\section{Task-conditioned observation policy}
\label{sec:agent}
\subsection{Episode state and evidence memory}
An episode begins with a question $q$, geographic ROI, initial state $s_0=(x_0,y_0,h_0)$, and fixed reference answer $y_q$. At step $t$, the renderer produces $I_t=\mathcal{R}(\mathcal{M},s_t)$; ChangeOS supplies predicted building locations and binary damage probabilities, and the hybrid answer pipeline using Qwen2.5-VL-7B-Instruct proposes an answer. The controller can answer, abstain, search, or request a dedicated recheck. The normalized terminal answer is scored against $y_q$, with abstentions counted as incorrect for the primary endpoint. Every policy uses the same frozen ChangeOS checkpoint and answer path.

Predicted objects are projected into geographic coordinates. For non-damage tasks, sightings are associated when their predicted centers are within 6 m or their geographic boxes have intersection-over-union at least 0.2; damage sightings use the public marked-building reference ID. Each track retains observation IDs and averages binary probabilities with inverse-GSD weights. Count and spatial answers use tracks detected in the current observation, so a stale historical detection cannot independently supply a current nearest target. The memory may still merge nearby buildings or miss objects, even when geometric ROI coverage is high.

\subsection{Task-conditioned recheck rule}
For a candidate recheck, T1\_TASK uses
\begin{equation}
 U_q=u_q(1-h_{t+1}/h_t)-0.05\,\widehat{T}_{\mathrm{move}}/60
 -(1-\widehat{C}_{\mathrm{ROI}}),
 \label{eq:taskutility}
\end{equation}
where $u_q\in[0,1]$ is a task-specific uncertainty score, $\widehat{T}_{\mathrm{move}}$ is configured-speed motion time in seconds, and $\widehat{C}_{\mathrm{ROI}}$ is predicted coverage after the move. A candidate executes only when $u_q\geq0.5$, $U_q\geq0.05$, and, for count or spatial questions, predicted ROI coverage is at least 0.98. Only descent contributes positively to the utility. The rule therefore triggers descent subject to task uncertainty and coverage constraints; it neither assigns value to same-height viewpoint changes nor estimates calibrated accuracy gains.

The uncertainty score is defined separately for each task. For damage, it is normalized binary entropy for the marked predicted building; a visible marked location with no matched detection receives $u_q=1$. For count, a capped Poisson-binomial distribution from detected buildings' damage probabilities gives entropy over $0,1,2,3+$. For a matched spatial candidate, the score is the maximum of its binary entropy and a GSD-dependent estimate of proximity to an eight-sector bearing boundary. Presence is answered from the wide view when possible; missing evidence enters search rather than a local recheck. Count uncertainty is conditional on detected objects and does not account for missed buildings. The spatial score uses a selected predicted target and does not fully represent uncertainty about which damaged building is nearest.

Damage rechecks may move toward the marked building by at most 40 m horizontally and descend. Count and spatial rechecks keep the geographic center and descend only while the task ROI remains covered. Each active policy has a cap of two dedicated rechecks. The stored policy identifier A0\_HOLD is referred to here as \emph{A0\_NO\_RECHECK}: it has a zero dedicated-recheck cap but, like the other policies, may use up to six search steps. Search may move and render a new image. Thus A0\_NO\_RECHECK is not a fixed-view control, and the reported recheck count is not total movement, rendering, or model-call cost.

The primary comparators are A1\_RANDOM (a seeded 0.5 recheck trigger), A2\_ALWAYS (request a feasible recheck), and A3U\_RAW\_ENTROPY (uncalibrated binary-entropy trigger). A3\_ENTROPY and A4\_CONFORMAL are supplemental policies. All seven configurations share the same question set and search cap, but their executed actions and subsequent answer-call counts may differ. This comparison therefore measures whole-policy outcomes at their realized operating points.
